% Curs 12 - Pagina 1 (3-I-1995).
% Fragment pentru % Curs 12 - Pagina 1 (3-I-1995).
% Fragment pentru % Curs 12 - Pagina 1 (3-I-1995).
% Fragment pentru % Curs 12 - Pagina 1 (3-I-1995).
% Fragment pentru \input{courses/pages/curs12-pagina1-ro}.
% Transcriere din manuscris: pastram ordinea si notatiile.
% Fara preambul; la inceput de curs nou, pagina noua.

\clearpage
\subsection*{Curs 12 - Pagina 1}

\begin{center}
\fbox{\textbf{CURS 12}}
\end{center}
\begin{flushright}
\textit{3-I-95}
\end{flushright}

\medskip
\noindent\underline{Generalit\u{a}\c{t}i asupra monoizilor}

\bigskip
\noindent\fbox{T.}\quad Fie $(M,\cdot,e)$ un monoid, $\varnothing\neq S\subseteq M$,
\[
\langle S\rangle
=\{x_1\cdots x_n\mid x_i\in S,\ n\geq 1\}\cup\{e\}.
\]
\noindent
\begin{enumerate}
\item[(1)] $\langle S\rangle$ este submonoid al lui $M$, $\langle S\rangle\leq M$.
\item[(2)] $N\leq M$, $S\subseteq N\ \Longrightarrow\ \langle S\rangle\subseteq N$.
\end{enumerate}

\noindent $S\neq\varnothing$, $S\subseteq M$ se nume\c{s}te \underline{submonoid}
dac\u{a} $S$ este stabil\u{a} la opera\c{t}ie \c{s}i $e\in S$.

\bigskip
\noindent\underline{Dem.} Evident.

\bigskip
\noindent\underline{Obs.}\quad
$\langle S\rangle$ este cel mai mic submonoid al lui $M$ care con\c{t}ine $S$.

\medskip
\noindent $\langle S\rangle$ se nume\c{s}te submonoidul lui $M$ generat de $S$.

\bigskip
\noindent\underline{Def.}\quad
$M$ - monoid, $S\subseteq M$ se nume\c{s}te \underline{sistem de generatori}
pentru $M$ dac\u{a}
\[
\langle S\rangle=M.
\]

\bigskip
\noindent\underline{Ex.}\quad Pentru $(\mathbb{N}^{*},\cdot,1)$,
un sistem de generatori este mul\c{t}imea numerelor prime.

\noindent Un alt sistem este mul\c{t}imea numerelor prime reunit\u{a} cu $\{1\}$.
% In exemplul manuscrisului: „{nr. prime} sau {nr. prime}+{1}”.
% Ultimele semne de la capatul randului sunt greu lizibile; fara completari.

\bigskip
\noindent\underline{Def.}\quad
Monoidul $M$ se nume\c{s}te \underline{finit generat}
dac\u{a} sistemul de generatori este finit.

\medskip
\noindent\underline{Obs.}\quad $M=\langle M\rangle$.

\bigskip
\noindent $R$ - congruen\c{t}\u{a} pe $M$
(rela\c{t}ie de echivalen\c{t}\u{a} compatibil\u{a} cu $\cdot$).
\[
\frac{M}{R}=\{\hat a\mid a\in M\},
\qquad \hat a\,\hat b=\widehat{ab},
\qquad \left(\frac{M}{R},\cdot,\hat e\right)
\text{ - monoidul factor.}
\]

\bigskip
\noindent\underline{Ex.}\quad $(\mathbb{Z},\cdot,1)$, $R_n$ - congruen\c{t}a modulo $n$,
\[
\frac{\mathbb{Z}}{R_n}=(\mathbb{Z}_n,\cdot,\hat 1).
\]
% Manuscrisul trateaza egalitatea de mai sus ca identificare cu monoidul factor.

\bigskip
\noindent\underline{Def.}\quad
Fie $(M,\cdot,e)$, $(M',\cdot,e')$ doi monoizi.

\noindent O aplica\c{t}ie $f:M\to M'$ se nume\c{s}te
\underline{morfism de monoizi} dac\u{a}:
\begin{enumerate}
\item[(1)] $f(xy)=f(x)f(y)$, $(\forall)\,x,y\in M$.
\item[(2)] $f(e)=e'$.
\end{enumerate}
% Definitia se incheie la sfarsitul paginii manuscrisului.
.
% Transcriere din manuscris: pastram ordinea si notatiile.
% Fara preambul; la inceput de curs nou, pagina noua.

\clearpage
\subsection*{Curs 12 - Pagina 1}

\begin{center}
\fbox{\textbf{CURS 12}}
\end{center}
\begin{flushright}
\textit{3-I-95}
\end{flushright}

\medskip
\noindent\underline{Generalit\u{a}\c{t}i asupra monoizilor}

\bigskip
\noindent\fbox{T.}\quad Fie $(M,\cdot,e)$ un monoid, $\varnothing\neq S\subseteq M$,
\[
\langle S\rangle
=\{x_1\cdots x_n\mid x_i\in S,\ n\geq 1\}\cup\{e\}.
\]
\noindent
\begin{enumerate}
\item[(1)] $\langle S\rangle$ este submonoid al lui $M$, $\langle S\rangle\leq M$.
\item[(2)] $N\leq M$, $S\subseteq N\ \Longrightarrow\ \langle S\rangle\subseteq N$.
\end{enumerate}

\noindent $S\neq\varnothing$, $S\subseteq M$ se nume\c{s}te \underline{submonoid}
dac\u{a} $S$ este stabil\u{a} la opera\c{t}ie \c{s}i $e\in S$.

\bigskip
\noindent\underline{Dem.} Evident.

\bigskip
\noindent\underline{Obs.}\quad
$\langle S\rangle$ este cel mai mic submonoid al lui $M$ care con\c{t}ine $S$.

\medskip
\noindent $\langle S\rangle$ se nume\c{s}te submonoidul lui $M$ generat de $S$.

\bigskip
\noindent\underline{Def.}\quad
$M$ - monoid, $S\subseteq M$ se nume\c{s}te \underline{sistem de generatori}
pentru $M$ dac\u{a}
\[
\langle S\rangle=M.
\]

\bigskip
\noindent\underline{Ex.}\quad Pentru $(\mathbb{N}^{*},\cdot,1)$,
un sistem de generatori este mul\c{t}imea numerelor prime.

\noindent Un alt sistem este mul\c{t}imea numerelor prime reunit\u{a} cu $\{1\}$.
% In exemplul manuscrisului: „{nr. prime} sau {nr. prime}+{1}”.
% Ultimele semne de la capatul randului sunt greu lizibile; fara completari.

\bigskip
\noindent\underline{Def.}\quad
Monoidul $M$ se nume\c{s}te \underline{finit generat}
dac\u{a} sistemul de generatori este finit.

\medskip
\noindent\underline{Obs.}\quad $M=\langle M\rangle$.

\bigskip
\noindent $R$ - congruen\c{t}\u{a} pe $M$
(rela\c{t}ie de echivalen\c{t}\u{a} compatibil\u{a} cu $\cdot$).
\[
\frac{M}{R}=\{\hat a\mid a\in M\},
\qquad \hat a\,\hat b=\widehat{ab},
\qquad \left(\frac{M}{R},\cdot,\hat e\right)
\text{ - monoidul factor.}
\]

\bigskip
\noindent\underline{Ex.}\quad $(\mathbb{Z},\cdot,1)$, $R_n$ - congruen\c{t}a modulo $n$,
\[
\frac{\mathbb{Z}}{R_n}=(\mathbb{Z}_n,\cdot,\hat 1).
\]
% Manuscrisul trateaza egalitatea de mai sus ca identificare cu monoidul factor.

\bigskip
\noindent\underline{Def.}\quad
Fie $(M,\cdot,e)$, $(M',\cdot,e')$ doi monoizi.

\noindent O aplica\c{t}ie $f:M\to M'$ se nume\c{s}te
\underline{morfism de monoizi} dac\u{a}:
\begin{enumerate}
\item[(1)] $f(xy)=f(x)f(y)$, $(\forall)\,x,y\in M$.
\item[(2)] $f(e)=e'$.
\end{enumerate}
% Definitia se incheie la sfarsitul paginii manuscrisului.
.
% Transcriere din manuscris: pastram ordinea si notatiile.
% Fara preambul; la inceput de curs nou, pagina noua.

\clearpage
\subsection*{Curs 12 - Pagina 1}

\begin{center}
\fbox{\textbf{CURS 12}}
\end{center}
\begin{flushright}
\textit{3-I-95}
\end{flushright}

\medskip
\noindent\underline{Generalit\u{a}\c{t}i asupra monoizilor}

\bigskip
\noindent\fbox{T.}\quad Fie $(M,\cdot,e)$ un monoid, $\varnothing\neq S\subseteq M$,
\[
\langle S\rangle
=\{x_1\cdots x_n\mid x_i\in S,\ n\geq 1\}\cup\{e\}.
\]
\noindent
\begin{enumerate}
\item[(1)] $\langle S\rangle$ este submonoid al lui $M$, $\langle S\rangle\leq M$.
\item[(2)] $N\leq M$, $S\subseteq N\ \Longrightarrow\ \langle S\rangle\subseteq N$.
\end{enumerate}

\noindent $S\neq\varnothing$, $S\subseteq M$ se nume\c{s}te \underline{submonoid}
dac\u{a} $S$ este stabil\u{a} la opera\c{t}ie \c{s}i $e\in S$.

\bigskip
\noindent\underline{Dem.} Evident.

\bigskip
\noindent\underline{Obs.}\quad
$\langle S\rangle$ este cel mai mic submonoid al lui $M$ care con\c{t}ine $S$.

\medskip
\noindent $\langle S\rangle$ se nume\c{s}te submonoidul lui $M$ generat de $S$.

\bigskip
\noindent\underline{Def.}\quad
$M$ - monoid, $S\subseteq M$ se nume\c{s}te \underline{sistem de generatori}
pentru $M$ dac\u{a}
\[
\langle S\rangle=M.
\]

\bigskip
\noindent\underline{Ex.}\quad Pentru $(\mathbb{N}^{*},\cdot,1)$,
un sistem de generatori este mul\c{t}imea numerelor prime.

\noindent Un alt sistem este mul\c{t}imea numerelor prime reunit\u{a} cu $\{1\}$.
% In exemplul manuscrisului: „{nr. prime} sau {nr. prime}+{1}”.
% Ultimele semne de la capatul randului sunt greu lizibile; fara completari.

\bigskip
\noindent\underline{Def.}\quad
Monoidul $M$ se nume\c{s}te \underline{finit generat}
dac\u{a} sistemul de generatori este finit.

\medskip
\noindent\underline{Obs.}\quad $M=\langle M\rangle$.

\bigskip
\noindent $R$ - congruen\c{t}\u{a} pe $M$
(rela\c{t}ie de echivalen\c{t}\u{a} compatibil\u{a} cu $\cdot$).
\[
\frac{M}{R}=\{\hat a\mid a\in M\},
\qquad \hat a\,\hat b=\widehat{ab},
\qquad \left(\frac{M}{R},\cdot,\hat e\right)
\text{ - monoidul factor.}
\]

\bigskip
\noindent\underline{Ex.}\quad $(\mathbb{Z},\cdot,1)$, $R_n$ - congruen\c{t}a modulo $n$,
\[
\frac{\mathbb{Z}}{R_n}=(\mathbb{Z}_n,\cdot,\hat 1).
\]
% Manuscrisul trateaza egalitatea de mai sus ca identificare cu monoidul factor.

\bigskip
\noindent\underline{Def.}\quad
Fie $(M,\cdot,e)$, $(M',\cdot,e')$ doi monoizi.

\noindent O aplica\c{t}ie $f:M\to M'$ se nume\c{s}te
\underline{morfism de monoizi} dac\u{a}:
\begin{enumerate}
\item[(1)] $f(xy)=f(x)f(y)$, $(\forall)\,x,y\in M$.
\item[(2)] $f(e)=e'$.
\end{enumerate}
% Definitia se incheie la sfarsitul paginii manuscrisului.
.
% Transcriere din manuscris: pastram ordinea si notatiile.
% Fara preambul; la inceput de curs nou, pagina noua.

\clearpage
\subsection*{Curs 12 - Pagina 1}

\begin{center}
\fbox{\textbf{CURS 12}}
\end{center}
\begin{flushright}
\textit{3-I-95}
\end{flushright}

\medskip
\noindent\underline{Generalit\u{a}\c{t}i asupra monoizilor}

\bigskip
\noindent\fbox{T.}\quad Fie $(M,\cdot,e)$ un monoid, $\varnothing\neq S\subseteq M$,
\[
\langle S\rangle
=\{x_1\cdots x_n\mid x_i\in S,\ n\geq 1\}\cup\{e\}.
\]
\noindent
\begin{enumerate}
\item[(1)] $\langle S\rangle$ este submonoid al lui $M$, $\langle S\rangle\leq M$.
\item[(2)] $N\leq M$, $S\subseteq N\ \Longrightarrow\ \langle S\rangle\subseteq N$.
\end{enumerate}

\noindent $S\neq\varnothing$, $S\subseteq M$ se nume\c{s}te \underline{submonoid}
dac\u{a} $S$ este stabil\u{a} la opera\c{t}ie \c{s}i $e\in S$.

\bigskip
\noindent\underline{Dem.} Evident.

\bigskip
\noindent\underline{Obs.}\quad
$\langle S\rangle$ este cel mai mic submonoid al lui $M$ care con\c{t}ine $S$.

\medskip
\noindent $\langle S\rangle$ se nume\c{s}te submonoidul lui $M$ generat de $S$.

\bigskip
\noindent\underline{Def.}\quad
$M$ - monoid, $S\subseteq M$ se nume\c{s}te \underline{sistem de generatori}
pentru $M$ dac\u{a}
\[
\langle S\rangle=M.
\]

\bigskip
\noindent\underline{Ex.}\quad Pentru $(\mathbb{N}^{*},\cdot,1)$,
un sistem de generatori este mul\c{t}imea numerelor prime.

\noindent Un alt sistem este mul\c{t}imea numerelor prime reunit\u{a} cu $\{1\}$.
% In exemplul manuscrisului: „{nr. prime} sau {nr. prime}+{1}”.
% Ultimele semne de la capatul randului sunt greu lizibile; fara completari.

\bigskip
\noindent\underline{Def.}\quad
Monoidul $M$ se nume\c{s}te \underline{finit generat}
dac\u{a} sistemul de generatori este finit.

\medskip
\noindent\underline{Obs.}\quad $M=\langle M\rangle$.

\bigskip
\noindent $R$ - congruen\c{t}\u{a} pe $M$
(rela\c{t}ie de echivalen\c{t}\u{a} compatibil\u{a} cu $\cdot$).
\[
\frac{M}{R}=\{\hat a\mid a\in M\},
\qquad \hat a\,\hat b=\widehat{ab},
\qquad \left(\frac{M}{R},\cdot,\hat e\right)
\text{ - monoidul factor.}
\]

\bigskip
\noindent\underline{Ex.}\quad $(\mathbb{Z},\cdot,1)$, $R_n$ - congruen\c{t}a modulo $n$,
\[
\frac{\mathbb{Z}}{R_n}=(\mathbb{Z}_n,\cdot,\hat 1).
\]
% Manuscrisul trateaza egalitatea de mai sus ca identificare cu monoidul factor.

\bigskip
\noindent\underline{Def.}\quad
Fie $(M,\cdot,e)$, $(M',\cdot,e')$ doi monoizi.

\noindent O aplica\c{t}ie $f:M\to M'$ se nume\c{s}te
\underline{morfism de monoizi} dac\u{a}:
\begin{enumerate}
\item[(1)] $f(xy)=f(x)f(y)$, $(\forall)\,x,y\in M$.
\item[(2)] $f(e)=e'$.
\end{enumerate}
% Definitia se incheie la sfarsitul paginii manuscrisului.
